% Intensité lumineuse derrière deux fentes de largeur non négligeable (corrigé du défi D22)
\begin{tikzpicture}
\begin{axis}[otMesures,otAxesBout,width=110mm,height=54mm,xmin=-16,xmax=16,ymin=0,ymax=1.12,
  xtick={-15,-10,...,15},ytick={0,0.5,1},yticklabels={0,0.5,1},
  xlabel={$x$ (mm)},ylabel={intensité relative}]
\addplot[otMuted,line width=0.6pt,dash pattern=on 2pt off 1.5pt,domain=-16:16,samples=401]
  {(abs(x)<0.001) ? 1 : (sin(deg(pi*x/12.66))/(pi*x/12.66))^2};
\addplot[otRed!80!black,line width=0.9pt,domain=-16:16,samples=801]
  {((abs(x)<0.001) ? 1 : (sin(deg(pi*x/12.66))/(pi*x/12.66))^2)*(cos(deg(pi*x/2.532)))^2};
\draw[otNavy,line width=0.5pt,{Stealth[length=1.3mm]}-{Stealth[length=1.3mm]}] (axis cs:-12.66,1.05) -- (axis cs:12.66,1.05);
\node[pcLegende,anchor=south,text=otNavy] at (axis cs:0,1.05) {tache centrale de diffraction : 9 franges brillantes};
\node[pcLegende,anchor=west,text=otMuted] at (axis cs:7.5,0.45) {enveloppe de diffraction};
\end{axis}
\end{tikzpicture}
\par\smallskip\legende{Les franges d'interférences sont modulées par la figure de diffraction de chaque fente ; la frange $k=\pm5$ tombe sur un zéro de diffraction et disparaît.}
